% 第1章 电路元件与电路基本定律（章节子文件，需被主文件调用）
\chapter{电路元件与电路基本定律}

\section{电流、电压与电功率}
\subsection{电流}
荷电质点的有序运动形成电流。为了衡量电流的大小，在时间$dt$内，通过某截面的电荷量的代数和为$dq$，则电流定义为
\[
i=\frac{dq}{dt}
\]
其方向规定为正电荷运动的方向。

电荷的$\mathrm{SI}$(\text{国际单位制})单位为库仑(\text{符号为}$\mathrm{C}$)，时间的$\mathrm{SI}$单位为秒时，电流的$\mathrm{SI}$单位为安培(\text{符号为}$\mathrm{A}$)，$1\mathrm{A}=1\mathrm{C}/1\mathrm{s}$。

电流的方向既可用箭头“$\rightarrow$”表示，也可用双下标表示。例如电流可从a点流向b点或相反，分别用$i_{ab}$和$i_{ba}$表示。在分析复杂电路时，难以判断电流的真实方向，而且有些电路中电流的真实方向随时间不断改变。为此，在列写电路方程时，可任意假设一个电流的方向，称为电流的参考方向，并用箭头标示在图中。选定参考方向后，电流便是代数量。电流的真实方向由电流代数量的符号和参考方向来判断。当按参考方向求得的电流为正值时，表明其真实方向与参考方向一致；否则相反。

量值和方向不随时间变化的电流称为直流。

\subsection{电压}
单位正电荷从电路中的一点移动到另一点，电场力对电荷做功。设电场力做功为$dW_{ab}$，则电压定义为
\[
u_{ab}=\frac{dW_{ab}}{dq}
\]
在式中，$u_{ab}$表示$\text{a、b}$两点间的电压，其物理意义是单位正电荷从$\text{a}$点移动到$\text{b}$点时电场力所做的功。

电压又称电位差，与电压相关的另一个物理量是电位。电位的物理意义是：在电路中可任选一点为参考点(\text{零电位点，常用符号}$\perp$\text{表示})，电路中某点与参考点之间的电压称为该点的电位，用$\phi$表示。

电压的方向规定为从高电位指向低电位的方向。电压参考方向可用双下标表示（如$u_{ab}$表示$\text{a}$点为高电位、$\text{b}$点为低电位）。

若电压的参考方向与真实方向一致，电压为正值；反之则为负值。在分析电路问题时，同样需要先设定电压的参考方向，再通过计算结果的符号判断真实方向。

\subsection{电动势}
从能量的角度看，电源是将其他形式的能量转换为电能的装置。设在$dt$的时间内，一电源使正电荷$dq$从负极经电源内部移至正极所做的功$dW$，电源的电动势定义为
\[
e=\frac{dW}{dq}
\]
即电源电动势的物理意义是单位正电荷从电源负极移至正极时，电源非静电力所做的功。

\subsection{电功率}
电功率是用以衡量电能转换速率的物理量，单位时间内电场力所做的功即为电功率，则电功率定义为
\[
p=\frac{dW}{dt}
\]

当电能被电路吸收时，功率为正值；当电路发出电能时，功率为负值。设某部分电路的电压参考方向与电流参考方向一致（关联参考方向），则该部分电路吸收的功率为$p=ui$；如果所选定的电压和电流的参考方向相反（非关联参考方向），则吸收的功率为$p=-ui$。

功率等于能量的转换速率，根据式可得在$t_0$到$t$时间内，电路吸收（电压、电流为关联参考方向时）或发出的能量为
\[
W(t)=\int_{t_0}^tp(\xi)d\xi=\int_{t_0}^tu(\xi)i(\xi)d\xi
\]
与判断功率正负类似，能量为正值时表示电路吸收能量，负值时表示发出能量。实际工程中，电气设备都有额定功率，如发电站、电动机等，若超过额定值，设备或部件可能因过热损坏。

\section{电阻元件}
在电路模型中，用电阻元件近似表示实际电阻器、灯泡、电炉等耗能器件的电气特性。在电路理论中，将流过相同电流的两个端子称为一个端口。线性二端电阻的端口电压与电流之间遵循欧姆定律，即二者成正比关系。

欧姆定律表达式为
\[
u=Ri
\]
或
\[
i=Gu
\]

若$u\text{、}i$参考方向相反，则欧姆定律应写为$u=-Ri$或$i=-Gu$。电阻与电导是反映同一电阻元件性能而互为倒数的两个参数，即
\[
RG=1
\]
线性电阻的端口特性（$u-i$曲线）是过原点的直线。

如果电阻$R=0$，则通过任意有限电流时其电压恒为零，故相当于短路，即短路线的电阻为零；如果$G=0$（$R\rightarrow\infty$），则施加任意有限电压时其电流恒为零，相当于开路。

电阻吸收的功率是$p=ui$，将欧姆定律代入，得
\[
p=ui=Ri^2=Gu^2
\]

如果$R\text{、}G$是正值，则$p$非负，这样的电阻消耗电功率，属于耗能元件。从能量角度看，直到$t$瞬间输入电阻中的全部电能将是
\[
W=\int_{-∞}^tp(\xi)d\xi=\int_{-∞}^tu(\xi)i(\xi)d\xi=R\int_{-∞}^ti^2(\xi)d\xi=G\int_{-∞}^tu^2(\xi)d\xi
\]
此式就是物理学中的焦耳定律，表明电阻消耗的电能全部转换为热能。

对于非负的电阻，在任一时刻，上式的结果非负。这表明：从全过程来看，这种元件只能吸收电能而不能发出电能。具有此种性质的元件，称为无源元件。

在电路模型中还存在一种负电阻，其$R<0$，能发出电能，属于有源元件（通常由电子电路实现）。

从广义上讲，只要端口电压、电流关系是代数方程，就属于电阻元件，其电压、电流方程一般形式为
\[
u=u(i)
\]
或
\[
i=i(u)
\]
上式中$u(i)$和$i(u)$分别表示电压是电流的函数、电流是电压的函数。如果电压、电流不是正比关系，即电压、电流关系不是过原点的直线，就称为非线性电阻。例如，半导体二极管就属于非线性电阻，其电压、电流关系为$i=I_0(e^{u/U_T}-1)$（$I_0$为反向饱和电流，$U_T$为热电压）。

\section{电容元件}
电容元件的电路符号中，固定电容和可变电容有不同表示。在电路模型中，用电容元件表示电容器的主要电气特性，即电荷与电压的关系。

当在电容器两端通过电源施加电压时，两极板上分别聚集起等量的正、负电荷，并在极板之间建立电场，储存有电场能量。当电源移去后，电荷和电场可以继续存在（理想电容），但实际电容器因漏电流会逐渐放电。电容能够储存电荷和电场能，是储能元件。

从物理学已知，当电容器填充线性介质（即介电常数与电场强度无关，如空气、橡皮、瓷、纱、胶木和云母等）时，它储存的电荷$q$与极板间电压$u$成正比。若电压的参考方向是从正极板指向负极板，则线性电容的$q-u$关系为
\[
q=Cu
\]
式中$C$是电容系数，其单位是$\mathrm{C/V}$(\text{库/伏})，称为法拉(\text{符号为}$\mathrm{F}$)，实际工程中常用微法（$\mu\mathrm{F}$，$1\mu\mathrm{F}=10^{-6}\mathrm{F}$）、皮法（$\mathrm{pF}$，$1\mathrm{pF}=10^{-12}\mathrm{F}$）。

平行板电容器的电容为
\[
C=\frac{\varepsilon A}{d}
\]
式中，$A$为每个极板的有效面积，$d$为极板间距离，$\varepsilon$为介质的介电常数（真空介电常数$\varepsilon_0=8.85\times10^{-12}\mathrm{F/m}$）。对线性介质，电容$C$是常量，称为线性电容；否则称为非线性电容。

线性电容的电荷-电压关系曲线是位于第$\mathrm{I}$、第$\mathrm{III}$象限且经过原点的直线。根据电流的定义，电流等于极板电荷的时间变化率
\[
i=\frac{dq}{dt}
\]
代入$q=Cu$，得
\[
i=C\frac{du}{dt}
\]
即线性电容的端口电流与端口电压的时间变化率成正比。当在端口上施加直流电压时，电压变化率为零，端口电流为零，电容等效成开路。

反过来，电荷$q$等于电流$i$对时间的积分，即
\[
q(t)=\int_{-∞}^ti(\xi)d\xi
\]
式表明$t$时刻电容上的电荷量是此时刻以前由电流充电和放电而“积累”的结果，所以电容具有记忆的性质，属于记忆元件。上式积分下限取$-∞$，意指从很久以前算起，电容上的电荷为零（初始状态）。

由$q=Cu$，得
\[
u(t)=\frac{q(t)}{C}=\frac{1}{C}\int_{-∞}^ti(\xi)d\xi
\]
式即为线性电容的$u-i$关系，式中$u\text{、}i$应为关联参考方向；否则要在式中添一负号（$u(t)=-\frac{1}{C}\int_{-∞}^ti(\xi)d\xi$）。

电阻元件是无记忆元件（电压仅由当前电流决定），而电容是记忆元件（电压与历史电流有关）。

如果研究$t_0$时刻以后的电容电压，可将积分拆分为两部分：
\[
u(t)=\frac{1}{C}\int_{-∞}^{t_0}i(\xi)d\xi+\frac{1}{C}\int_{t_0}^ti(\xi)d\xi=u(t_0)+\frac{1}{C}\int_{t_0}^ti(\xi)d\xi
\]
式中，$u(t_0)$表示$t_0$时刻以前全部电流对电压的影响，称为电容在$t_0$时的初始电压。

输入线性电容的功率为
\[
p=ui=Cu\frac{du}{dt}
\]
根据能量与功率的关系，由式可以求得截止到$t$时刻电容吸收的总能量
\[
W_e(t)=\int_{-∞}^tp(\xi)d\xi=\int_{-∞}^tCu(\xi)\frac{du(\xi)}{d\xi}d\xi=C\int_{u(-∞)}^{u(t)}u du
\]
在$t=-∞$时电容尚未充电，所以$u(-∞)=0$，故电容吸收的总能量为
\[
W_e(t)=\frac{1}{2}Cu^2=\frac{q^2}{2C}
\]
式就是电容储存电场能量的计算公式，表明电容储存的能量仅与当前电压有关，与充电过程无关。

\section{电感元件}
在电机、变压器、电磁铁等电气设备中，线圈是核心部件。尽管实际线圈存在电阻，但其主要电气特性是通过电流产生磁场、储存磁场能量。在电路模型中，用电感元件表示线圈的这一核心特性。

如果线圈的磁场存在于线性介质中（如空气、非铁磁材料），则磁链与电流成正比，称为线性电感；否则（如铁磁材料）称为非线性电感。

设线性电感的磁链$\Psi$与电流$i$的参考方向符合右手螺旋定则，则二者之间的关系为
\[
\Psi=Li
\]
式中$L$为电感系数，磁链的单位是韦伯(\text{符号}$\mathrm{Wb}$)，电流的单位是安培($\mathrm{A}$)时，电感的$\mathrm{SI}$单位是亨利(\text{符号}$\mathrm{H}$)，$1\mathrm{H}=1\mathrm{Wb/A}$，实际常用毫亨（$\mathrm{mH}$）、微亨（$\mu\mathrm{H}$）。

电感的特性还可以用图形表示，与式对应的磁链—电流关系是一条通过$\Psi-i$平面原点的直线且位于第$\mathrm{I}$、第$\mathrm{III}$象限。

根据电磁感应定律，磁链随时间的变化要在电感中产生感应电动势，记作$e$。当电压$u$、电动势$e$与电流$i$取关联参考方向（电压参考方向与电动势相反），则
\[
u=-e=-(-\frac{d\Psi}{dt})=\frac{d\Psi}{dt}
\]
对线性电感，代入$\Psi=Li$，得
\[
u=L\frac{di}{dt}
\]
即线性电感的端口电压与端口电流的时间变化率成正比。当电感中流过直流电流时，电流变化率为零，端口电压为零，电感等效成短路。

如果从电压求磁链或电流，则需积分，即
\[
\Psi(t)=\int_{-∞}^tu(\xi)d\xi=\Psi(t_0)+\int_{t_0}^tu(\xi)d\xi
\]
\[
i(t)=\frac{1}{L}\int_{-∞}^tu(\xi)d\xi=i(t_0)+\frac{1}{L}\int_{t_0}^tu(\xi)d\xi
\]
此两式表明，电感中某一瞬间的磁链和电流决定于此瞬间以前的全部电压，因此电感也属于记忆元件。

上述$u-i$关系要求$u\text{、}i$为关联参考方向；否则要在式中添一负号（$u=-L\frac{di}{dt}$）。

线性电感吸收的功率为
\[
p=ui=Li\frac{di}{dt}
\]
截止到$t$时刻电感吸收的能量为
\[
W_m = \int_{-\infty}^{t} p(\xi) \, d\xi = L \int_{i(-\infty)}^{i(t)} i(\xi) \, di(\xi) = \frac{1}{2} L i^2 \bigg|_{i(-\infty)}^{i(t)}
\]
在$t=-∞$时，电感电流为零，所以电感吸收的磁场能量为
\[
W_m=\frac{1}{2}Li^2=\frac{\Psi^2}{2L}
\]
式是电感储存磁场能量的计算公式。当电感电流的绝对值增大时，电感吸收能量；电流减小时，释放能量，所以电感元件也是一个储能元件（无源元件）。

\section{独立电源}
电源模型有电压源和电流源两种，均为独立电源（源电压/电流不受电路其他量控制）。

\subsection{电压源}
理想电压源的端口电压与电流大小无关，或为常量或随时间按确定规律变化。

在电路模型中，电压源的符号为“+”、“-”极表示，其端口特性是：无论流过的电流是多少，端口电压恒等于源电压$u_s$。

电压源能够提供确定的电源电压$u_s$，称为源电压。所谓“确定”是指源电压$u_s$与流过电压源的电流无关，电压源的电流将由与其相连的电路来确定。$u_s$可以是常量，称为直流电压源（或恒定电压源），记作$u_s=U_s$；也可以是时变量，记作$u_s=u_s(t)$（如正弦电压源）。源电压$u_s$的方向是从“+”极到“-”极。

电压源的端口特性可以用解析式表示
\[
u=u_s
\]
式中不含电流$i$，表示源电压不受电流的制约。源电压为常量时，端口电压与电流之间的关系可用平行于电流坐标轴的直线表示，称为电源的端口特性曲线。当源电压$u_s=0$即电压源置零时，电压源的作用相当于短路。

在关联参考方向下（电流从“+”极流入），电压源吸收的功率是$p=u_si$；若电流从“+”极流出（非关联参考方向），则$p=-u_si$为发出的功率。当$p>0$时，电压源吸收能量（充电）；$p<0$时，发出能量（供电）。

\subsection{电流源}
理想电流源的端口电流与端口电压无关，恒等于源电流$i_s$。

在电路模型中，电流源的符号为箭头表示电流方向，其端口特性是：无论端口电压是多少，端口电流恒等于源电流$i_s$。

电流源能够提供确定的端口电流$i_s$，称为源电流。这里“确定”是指$i_s$与电流源端口电压无关，电流源的端口电压决定于它所接的外电路。$i_s$可以是常量即恒定电流，记作$i_s=I_s$，也可以是时变量，记作$i_s=i_s(t)$。

电流源端口特性的解析表达式为
\[
i=i_s
\]
式不含电压$u$，表明电流源的端口电流不受端口电压的制约。如果源电流$i_s$为常量，则式对应$u-i$平面上一条平行于电压坐标轴的直线。当$i_s=0$即电流源置零时，电流源的作用相当于开路。

在关联参考方向下（电流从“+”极流入），电流源吸收的功率为$p=ui_s$；非关联参考方向下，发出的功率为$p=-ui_s$。当特性曲线位于第$\mathrm{I}$、第$\mathrm{III}$象限时，$p>0$，电流源处于充电状态；当特性曲线位于第$\mathrm{II}$、第$\mathrm{IV}$象限时，$p<0$，电流源处于供电状态。

电压源和电流源作为元件模型，是实际电源（如电池、发电机）的理想化近似：实际电压源存在内阻，端口电压随电流增大而降低；实际电流源存在内电导，端口电流随电压增大而减小。

\section{受控电源}
在电路模型中还要用到另一类电源，它们的源电压或源电流受电路中另一处的电压或电流控制，这类电源称为受控电源或非独立电源。

受控源是一种用来表示晶体管、运算放大器等电子器件的电路模型，本书只讨论线性受控源，即源电压（流）与控制电压（流）成正比关系。用菱形符号表示受控源，以便与独立源相区分。

线性受控源分为四种类型，其端口方程分别为：
\begin{enumerate}
    \item 电压控制电压源（\textbf{VCVS}）：$\begin{cases}u_s = \alpha u_c \\i_c = 0\end{cases}$（$\alpha$为电压放大系数，无量纲）
    \item 电流控制电压源（\textbf{CCVS}）：$\begin{cases}u_s = r i_c \\u_c = 0\end{cases}$（$r$为转移电阻，单位$\Omega$）
    \item 电压控制电流源（\textbf{VCCS}）：$\begin{cases}i_s = g u_c \\i_c = 0\end{cases}$（$g$为转移电导，单位$\mathrm{S}$）
    \item 电流控制电流源（\textbf{CCCS}）：$\begin{cases}i_s = \beta i_c \\u_c = 0\end{cases}$（$\beta$为电流放大系数，无量纲）
\end{enumerate}
式中带脚注$s$的变量表示源量，带脚注$c$的变量表示控制量。参量$\alpha\text{、}r\text{、}g\text{、}\beta$都是常量，由器件特性决定。

受控源是双端口元件：一个端口是控制端口（输入端口），另一个是受控端口（输出端口）。控制端口若为开路（$i_c=0$），则是电压控制型；若为短路（$u_c=0$），则是电流控制型。

从电路模型角度，受控源本身不消耗/提供能量（理想情况），其能量由控制电路提供，因此受控源属于无源元件（但由受控源组成的电路可能是有源的）。

\section{基尔霍夫定律}
1847年基尔霍夫建立了电路结构对电流、电压的约束关系，分别称为基尔霍夫电流定律（KCL）和基尔霍夫电压定律（KVL），统称基尔霍夫定律，是分析集中参数电路的基本定律。

\subsection{电路结构基本概念}
电路元件通过端子相互连接便形成具有特定结构的电路。本书将每个二端元件称为一条支路，若干支路的连接点称为节点（两个支路的连接点称为简单节点），由支路构成的闭合路径称为回路，内部不含其他支路的回路称为网孔。

一条支路的电流是指流过该支路的电流，支路的电压是指支路两个端子之间的电压。

\subsection{基尔霍夫电流定律（KCL）}
基尔霍夫电流定律表述为：在集中参数电路中，任一时刻流出（或流入）任一节点的支路电流代数和等于零，即
\[
\sum i_k=0
\]
通常规定：电流参考方向为流出节点时取正号，流入节点时取负号（也可反之，只需统一）。

式称为节点的$\mathrm{KCL}$方程，它反映了与任意一个节点相连的所有支路电流所满足的电荷守恒约束（节点不能储存电荷）。

$\mathrm{KCL}$还可推广用于流出（或流入）闭合边界的支路电流（称为“广义节点”），表述为：在集中参数电路中，任一时刻流出（或流入）任一闭合边界$S$的支路电流代数和等于零，即
\[
\sum_{k\in S} i_k=0
\]
式称为广义基尔霍夫电流定律。

如将流出节点的电流记为正、流入记为负，则$\mathrm{KCL}$也可表述为：任一时刻，流出任一节点（或闭合边界）的电流之和等于流入该节点（或闭合边界）的电流之和，即
\[
\sum i_{\text{流出}}=\sum i_{\text{流入}}
\]

\subsection{基尔霍夫电压定律（KVL）}
基尔霍夫电压定律表述为：在集中参数电路中，任一时刻沿任一回路各支路电压的代数和等于零，即
\[
\sum u_k=0
\]
通常规定：当支路电压的参考方向与回路绕行方向相同时，$u_k$前面取正号；反之取负号。

式称为回路的$\mathrm{KVL}$方程，它是任一回路所包含的支路电压必须满足的能量守恒约束（单位电荷沿回路绕行一周，电场力做功为零）。

推广到一般情况就是：沿任一回路，各支路电压降的代数和等于电压升的代数和，即
\[
\sum u_{\text{降}}=\sum u_{\text{升}}
\]
进一步推广：集中参数电路中，任意两点之间的电压具有确定值，与计算路径无关（电压的单值性）。

对电路中的每个网孔或回路均可列出$\mathrm{KVL}$方程，但并非所有方程都是独立的。可以证明：对于具有$b$条支路、$n$个节点的电路，独立$\mathrm{KVL}$方程的个数为$b-(n-1)$（平面电路的网孔数等于独立$\mathrm{KVL}$方程数）。

基尔霍夫定律仅与电路的拓扑结构有关，与元件的性质无关（线性/非线性、时变/时不变均适用）。